% Options for packages loaded elsewhere
\PassOptionsToPackage{unicode}{hyperref}
\PassOptionsToPackage{hyphens}{url}
\documentclass[
]{article}
\usepackage{xcolor}
\usepackage{amsmath,amssymb}
\setcounter{secnumdepth}{5}
\usepackage{iftex}
\ifPDFTeX
  \usepackage[T1]{fontenc}
  \usepackage[utf8]{inputenc}
  \usepackage{textcomp} % provide euro and other symbols
\else % if luatex or xetex
  \usepackage{unicode-math} % this also loads fontspec
  \defaultfontfeatures{Scale=MatchLowercase}
  \defaultfontfeatures[\rmfamily]{Ligatures=TeX,Scale=1}
\fi
\usepackage{lmodern}
\ifPDFTeX\else
  % xetex/luatex font selection
\fi
% Use upquote if available, for straight quotes in verbatim environments
\IfFileExists{upquote.sty}{\usepackage{upquote}}{}
\IfFileExists{microtype.sty}{% use microtype if available
  \usepackage[]{microtype}
  \UseMicrotypeSet[protrusion]{basicmath} % disable protrusion for tt fonts
}{}
\makeatletter
\@ifundefined{KOMAClassName}{% if non-KOMA class
  \IfFileExists{parskip.sty}{%
    \usepackage{parskip}
  }{% else
    \setlength{\parindent}{0pt}
    \setlength{\parskip}{6pt plus 2pt minus 1pt}}
}{% if KOMA class
  \KOMAoptions{parskip=half}}
\makeatother
\usepackage{longtable,booktabs,array}
\newcounter{none} % for unnumbered tables
\usepackage{calc} % for calculating minipage widths
% Correct order of tables after \paragraph or \subparagraph
\usepackage{etoolbox}
\makeatletter
\patchcmd\longtable{\par}{\if@noskipsec\mbox{}\fi\par}{}{}
\makeatother
% Allow footnotes in longtable head/foot
\IfFileExists{footnotehyper.sty}{\usepackage{footnotehyper}}{\usepackage{footnote}}
\makesavenoteenv{longtable}
\usepackage{graphicx}
\makeatletter
\newsavebox\pandoc@box
\newcommand*\pandocbounded[1]{% scales image to fit in text height/width
  \sbox\pandoc@box{#1}%
  \Gscale@div\@tempa{\textheight}{\dimexpr\ht\pandoc@box+\dp\pandoc@box\relax}%
  \Gscale@div\@tempb{\linewidth}{\wd\pandoc@box}%
  \ifdim\@tempb\p@<\@tempa\p@\let\@tempa\@tempb\fi% select the smaller of both
  \ifdim\@tempa\p@<\p@\scalebox{\@tempa}{\usebox\pandoc@box}%
  \else\usebox{\pandoc@box}%
  \fi%
}
% Set default figure placement to htbp
\def\fps@figure{htbp}
\makeatother
\setlength{\emergencystretch}{3em} % prevent overfull lines
\providecommand{\tightlist}{%
  \setlength{\itemsep}{0pt}\setlength{\parskip}{0pt}}
\usepackage{bookmark}
\IfFileExists{xurl.sty}{\usepackage{xurl}}{} % add URL line breaks if available
\urlstyle{same}
\hypersetup{
  hidelinks,
  pdfcreator={LaTeX via pandoc}}

\author{}
\date{}

\begin{document}

\includegraphics[width=4.2031in,height=0.94625in]{media/image1.png}

Universidad Internacional de La Rioja

Escuela Superior de Ingeniería y Tecnología

Máster Universitario en Ingeniería Matemática y Computación

Sistema híbrido de detección de anomalías y Momentum Trading con
algoritmos de clustering (K-Means/DBSCAN)

{\def\LTcaptype{none} % do not increment counter
\begin{longtable}[]{@{}
  >{\raggedright\arraybackslash}p{(\linewidth - 2\tabcolsep) * \real{0.4532}}
  >{\raggedright\arraybackslash}p{(\linewidth - 2\tabcolsep) * \real{0.4864}}@{}}
\toprule\noalign{}
\begin{minipage}[b]{\linewidth}\raggedright
Trabajo fin de estudio presentado por:
\end{minipage} & \begin{minipage}[b]{\linewidth}\raggedright
Espejo Cuenca, Marcel

Gómez Hernández, Raúl

Martínez Díez, Diego
\end{minipage} \\
\midrule\noalign{}
\endhead
\bottomrule\noalign{}
\endlastfoot
Tipo de trabajo: & Desarrollo Práctico \\
Director/a: & Martínez Cerdá, Juan Francisco \\
Fecha: & 15/04/2026 \\
\end{longtable}
}

Índice de contenidos

\hyperref[contexto-y-estado-del-arte]{1. Contexto y estado del arte
\hyperref[contexto-y-estado-del-arte]{7}}

\hyperref[el-momentum-en-los-mercados-financieros-y-el-riesgo-de-anomaluxedas]{1.1.
El Momentum en los mercados financieros y el riesgo de anomalías
\hyperref[el-momentum-en-los-mercados-financieros-y-el-riesgo-de-anomaluxedas]{7}}

\hyperref[fundamentaciuxf3n-matemuxe1tica-del-momentum]{1.1.1.
Fundamentación matemática del Momentum
\hyperref[fundamentaciuxf3n-matemuxe1tica-del-momentum]{7}}

\hyperref[asimetruxeda-de-riesgo-y-el-concepto-de-momentum-crashes]{1.1.2.
Asimetría de riesgo y el concepto de Momentum Crashes
\hyperref[asimetruxeda-de-riesgo-y-el-concepto-de-momentum-crashes]{9}}

\hyperref[ruido-estructural-anomaluxedas-y-la-necesidad-de-filtrado-topoluxf3gico]{1.1.3.
Ruido estructural, anomalías y la necesidad de filtrado topológico
\hyperref[ruido-estructural-anomaluxedas-y-la-necesidad-de-filtrado-topoluxf3gico]{11}}

\hyperref[fundamentos-de-machine-learning-no-supervisado]{1.2.
Fundamentos de Machine Learning no supervisado
\hyperref[fundamentos-de-machine-learning-no-supervisado]{12}}

\hyperref[algoritmo-de-agrupamiento-k-means]{1.2.1. Algoritmo de
Agrupamiento K-Means \hyperref[algoritmo-de-agrupamiento-k-means]{12}}

\hyperref[filtro-de-densidad-topoluxf3gica-dbscan]{1.2.2. Filtro de
densidad topológica DBSCAN
\hyperref[filtro-de-densidad-topoluxf3gica-dbscan]{14}}

\hyperref[estado-del-arte-y-justificaciuxf3n-de-la-arquitectura-propuesta]{1.3.
Estado del arte y justificación de la arquitectura propuesta
\hyperref[estado-del-arte-y-justificaciuxf3n-de-la-arquitectura-propuesta]{15}}

\hyperref[el-problema-del-momentum-tradicional-en-entornos-voluxe1tiles]{1.3.1.
El problema del Momentum tradicional en entornos volátiles
\hyperref[el-problema-del-momentum-tradicional-en-entornos-voluxe1tiles]{16}}

\hyperref[el-intento-de-soluciuxf3n-mediante-k-means-y-sus-limitaciones]{1.3.2.
El intento de solución mediante K-Means y sus limitaciones
\hyperref[el-intento-de-soluciuxf3n-mediante-k-means-y-sus-limitaciones]{16}}

\hyperref[propuesta-de-soluciuxf3n-arquitectura-huxedbrida-y-marco-de-evaluaciuxf3n]{1.3.3.
Propuesta de Solución: Arquitectura Híbrida y Marco de Evaluación
\hyperref[propuesta-de-soluciuxf3n-arquitectura-huxedbrida-y-marco-de-evaluaciuxf3n]{17}}

\hyperref[objetivos-concretos-y-metodologuxeda-de-trabajo]{2. Objetivos
concretos y metodología de trabajo
\hyperref[objetivos-concretos-y-metodologuxeda-de-trabajo]{18}}

\hyperref[objetivo-general]{2.1. Objetivo general
\hyperref[objetivo-general]{18}}

\hyperref[objetivos-especuxedficos]{2.2. Objetivos específicos
\hyperref[objetivos-especuxedficos]{18}}

\hyperref[metodologuxeda-del-trabajo]{2.3. Metodología del trabajo
\hyperref[metodologuxeda-del-trabajo]{19}}

\hyperref[obtenciuxf3n-de-datos]{2.3.1. Obtención de datos
\hyperref[obtenciuxf3n-de-datos]{19}}

\hyperref[procesamiento-y-preparaciuxf3n-de-datos]{2.3.2. Procesamiento
y preparación de datos
\hyperref[procesamiento-y-preparaciuxf3n-de-datos]{20}}

\hyperref[anuxe1lisis-exploratorio-de-datos-eda]{2.3.3. Análisis
exploratorio de datos (EDA)
\hyperref[anuxe1lisis-exploratorio-de-datos-eda]{20}}

\hyperref[implementaciuxf3n-de-la-arquitectura-huxedbrida]{2.3.4.
Implementación de la arquitectura híbrida
\hyperref[implementaciuxf3n-de-la-arquitectura-huxedbrida]{20}}

\hyperref[estrategia-de-momentum-y-detecciuxf3n-de-anomaluxedas]{2.3.5.
Estrategia de momentum y detección de anomalías
\hyperref[estrategia-de-momentum-y-detecciuxf3n-de-anomaluxedas]{21}}

\hyperref[evaluaciuxf3n-del-sistema]{2.3.6. Evaluación del sistema
\hyperref[evaluaciuxf3n-del-sistema]{21}}

\hyperref[entorno-de-desarrollo-y-herramientas]{2.3.7. Entorno de
desarrollo y herramientas
\hyperref[entorno-de-desarrollo-y-herramientas]{21}}

\hyperref[_Toc227174287]{3. Referencias bibliográficas
\hyperref[_Toc227174287]{22}}

Índice de figuras

\hyperref[_Toc227174288]{Figura 1. Evolución de las caídas catastróficas
en estrategias de momentum. El desplome en (Panel A) se explica por el
rebote explosivo de los activos con peor rendimiento previo mostrado en
(Panel B). Fuente: Daniel y Moskowitz (2016).
\hyperref[_Toc227174288]{10}}

\hyperref[_Toc227174289]{Figura 2: Clasificación de activos en el
espacio riesgo-rentabilidad mediante el algoritmo KMeans. Los puntos
sólidos representan los centroides µi calculados según la Ecuación 2,
mientras que las regiones sombreadas delimitan las fronteras de decisión
de cada clúster. Fuente: Elaboración propia.
\hyperref[_Toc227174289]{13}}

\hyperref[_Toc227174290]{Figura 3: Clasificación topológica en DBSCAN:
Identificación de puntos núcleo, frontera y ruido en función del radio ϵ
y MinPts. Fuente: Elaboración propia. \hyperref[_Toc227174290]{15}}

Índice de tablas

\hyperref[_Toc227174291]{Tabla 1 Rendimientos de las carteras
superpuestas J\textbackslash/K de momentum. Fuente: Jegadeesh y Titman
(1993). \hyperref[_Toc227174291]{8}}

\textbf{\hfill\break
}

Índice de ecuaciones

\hyperref[_Toc227174292]{Ecuación 1. Minimización de la suma de los
errores cuadráticos \hyperref[_Toc227174292]{12}}

\hyperref[_Ref227165937]{Ecuación 2. Media aritmética de los puntos del
clúster \hyperref[_Ref227165937]{12}}

\hyperref[_Toc227174294]{Ecuación 3. Distancia euclídea
\hyperref[_Toc227174294]{13}}

\hyperref[_Toc227174295]{Ecuación 4. Vecindad
\hyperref[_Toc227174295]{14}}

\section{Contexto y estado del arte}\label{contexto-y-estado-del-arte}

El presente capítulo trata los fundamentos teóricos y matemáticos que
sustentan la investigación, así como la revisión de la literatura
contemporánea. En primer lugar, se analiza la naturaleza del
\emph{momentum} y su vulnerabilidad ante regímenes de alta volatilidad.
Posteriormente, se muestran las herramientas de \emph{Machine Learning}
no supervisado propuestas (K-Means y DBSCAN), para finalizar con un
análisis del estado del arte que justifica la arquitectura algorítmica
híbrida de este trabajo.

\subsection{El Momentum en los mercados financieros y el riesgo de
anomalías}\label{el-momentum-en-los-mercados-financieros-y-el-riesgo-de-anomaluxedas}

La persistencia de los rendimientos en los mercados de renta variable ha
sido objeto de profundo debate en la literatura financiera. La Hipótesis
del Mercado Eficiente (EMH) establece que los precios reflejan
instantáneamente toda la información disponible, imposibilitando la
obtención sistemática de rendimientos anormales. Sin embargo, la
evidencia empírica ha demostrado que los inversores tienden a reaccionar
de forma ineficiente (por sobrerreacción o infrarreacción) ante la nueva
información, generando tendencias explotables matemáticamente. A la
captura sistemática de estas tendencias se le conoce como inversión de
\emph{momentum}.

\subsubsection{\texorpdfstring{Fundamentación matemática del Momentum
}{Fundamentación matemática del Momentum }}\label{fundamentaciuxf3n-matemuxe1tica-del-momentum}

El \emph{momentum} se fundamenta en la premisa de que los activos que
han obtenido un rendimiento superior en el pasado reciente continuarán
haciéndolo en el futuro cercano, mientras que aquellos con un
rendimiento deficiente seguirán una trayectoria bajista. La
formalización académica de esta estrategia fue establecida de forma
pionera por (Jegadeesh, 1993), quienes demostraron que comprar acciones
ganadoras y vender acciones perdedoras genera rendimientos anormales
positivos y estadísticamente significativos.

La arquitectura matemática del modelo clásico de (Jegadeesh, 1993) se
define mediante una estrategia de \(J\) meses de formación y \(K\) meses
de mantenimiento (estrategia \(J\text{/}K\) ) . En el instante \(t\), el
universo de activos se clasifica en orden ascendente en función de sus
rendimientos acumulados durante los pasados \(J\) meses. Basándose en
este ranking, se forman diez carteras (deciles) equiponderadas. El decil
superior se define como la cartera de "Ganadores" (Winners) y el decil
inferior como la cartera de "Perdedores" (Losers).

La estrategia consiste en construir una cartera de inversión de coste
cero (autofinanciada), tomando una posición larga en el decil de
ganadores y una posición corta en el decil de perdedores, mantenida
durante \(K\) periodos. Para maximizar el poder estadístico y evitar el
sesgo de un único punto de entrada, se propone el uso de carteras
superpuestas (overlapping portfolios). En cualquier mes \(t\), la
estrategia mantiene una serie de carteras seleccionadas en el mes actual
y en los \(K - 1\ \)meses anteriores, rebalanceando mensualmente una
fracción \(1\text{/}K\) de los pesos del hiperespacio de activos.

\begin{figure}
\centering
\includegraphics[width=6.29931in,height=4.74653in]{media/image2.PNG}
\caption{\protect\phantomsection\label{_Toc227174291}{}Tabla
Rendimientos de las carteras superpuestas J\textbackslash/K de momentum.
Fuente: Jegadeesh y Titman (1993).}
\end{figure}

Esta formulación lineal clásica ha demostrado ser rentable, pero asume
implícitamente una distribución estacionaria del riesgo, una suposición
que resulta matemáticamente inestable en la práctica.

\subsubsection{Asimetría de riesgo y el concepto de Momentum
Crashes}\label{asimetruxeda-de-riesgo-y-el-concepto-de-momentum-crashes}

A pesar de sus altos rendimientos promedio y ratios de Sharpe
atractivos, la estrategia de momentum presenta un riesgo de cola. La
distribución de sus rendimientos está fuertemente sesgada hacia la
izquierda, lo que significa que el modelo experimenta pérdidas
infrecuentes pero devastadoras, conocidas en la literatura como Momentum
Crashes (Daniel, 2016).

Daniel y Moskowitz (Daniel, 2016) demuestran que estos colapsos no son
eventos aleatorios (saltos de Poisson), sino que son parcialmente
predecibles al concentrarse en "estados de pánico" del mercado; es
decir, tras caídas prolongadas (bear markets) combinadas con picos
extremos de volatilidad esperada. La anatomía de un crash tiene una
forma contraintuitiva: el colapso de la estrategia no ocurre cuando el
mercado sigue cayendo, sino exactamente cuando rebota de forma abrupta.

La justificación matemática radica en la variación dinámica de la beta (
\(\beta\) ) de las carteras extremas. Tras un declive severo, las
acciones del decil de "Perdedores" han sufrido grandes caídas, quedando
a menudo en situaciones de estrés financiero. En este punto, el capital
de estas empresas actúa matemáticamente como una call option (opción de
compra) fuera del dinero sobre los activos de la empresa. Mediante la
introducción de variables indicadoras de mercado bajista ( \(I_{B}\) ) y
alcista ( \(I_{U}\) ), Daniel y Moskowitz (Daniel, 2016) demuestran que
en estados de pánico, la cartera de perdedores presenta una asimetría
extrema: su beta a la baja es moderada, pero su beta al alza se dispara
ante cualquier rebote positivo.

Dado que el algoritmo de momentum obliga a mantener una posición corta
sobre estos perdedores, la estrategia se encuentra efectivamente "corta
en una opción de compra" sobre el índice del mercado. Cuando ocurre un
rebote (ej. marzo-mayo de 2009), la cartera de perdedores sube
destruyendo el margen de la posición corta, mientras que la cartera de
ganadores apenas sube, generando un elevado drawdown en la estrategia.

\includegraphics[width=6.29931in,height=3.76736in]{media/image3.PNG}\includegraphics[width=6.29931in,height=3.70417in]{media/image4.PNG}

\protect\phantomsection\label{_Toc227174288}{}Figura . Evolución de las
caídas catastróficas en estrategias de momentum. El desplome en (Panel
A) se explica por el rebote explosivo de los activos con peor
rendimiento previo mostrado en (Panel B). Fuente: Daniel y Moskowitz
(2016).

\subsubsection{Ruido estructural, anomalías y la necesidad de filtrado
topológico}\label{ruido-estructural-anomaluxedas-y-la-necesidad-de-filtrado-topoluxf3gico}

La evidencia de los \emph{Momentum Crashes} y la inestabilidad de los
modelos lineales subraya un desafío en la econometría financiera: el
bajo \emph{signal-to-noise ratio} de las series temporales. Como expone
(López de Prado, 2018), los datos financieros están contaminados de
forma continua por eventos de microestructura y anomalías
idiosincrásicas que no representan la dinámica real del mercado.

En este contexto, es crucial distinguir entre un cambio de régimen
sistémico (donde la mayoría de los activos se desplazan conjuntamente,
formando estructuras densas) y el ruido estructural o valores atípicos
individuales (\emph{outliers}). Los algoritmos de partición
tradicionales, como K-Means, presentan una alta sensibilidad a estos
últimos; la presencia de una sola acción con una rentabilidad extrema e
injustificada puede desplazar significativamente los centroides,
distorsionando la clasificación de todo el universo de activos.

Para solucionar esta limitación, la literatura técnica más reciente
propone el uso de algoritmos basados en densidad. Específicamente, el
algoritmo DBSCAN (\emph{Density-Based Spatial Clustering of Applications
with Noise}) permite una aproximación topológica al filtrado de datos. A
diferencia de los métodos de exclusión manuales o estáticos propuestos
por autores como (Aslam, 2025), DBSCAN automatiza la identificación de
anomalías basándose en la proximidad espacial de las observaciones
(Shiraj, y otros, 2024). Aquellos activos que presenten movimientos de
precio extremos, \emph{flash crashes} o patrones de volatilidad anómalos
que no cuenten con una densidad mínima de ``vecinos'' con comportamiento
similar, son etiquetados matemáticamente como ruido (\emph{noise}).

Por consiguiente, se justifica la implementación de una arquitectura
híbrida en este Trabajo. El uso de DBSCAN actúa como un paso de
preprocesamiento para purgar las anomalías individuales y ``cisnes
negros'' puntuales. Una vez eliminada esta contaminación estadística, el
algoritmo K-Means puede clasificar los activos en regímenes de mercado
robustos, garantizando que los centroides representen tendencias reales
y no se vean afectados por el ruido estructural del mercado.

\subsection{Fundamentos de Machine Learning no
supervisado}\label{fundamentos-de-machine-learning-no-supervisado}

Para abordar la naturaleza no lineal y ruidosa de los mercados
financieros descrita en la sección anterior, esta investigación propone
un sistema híbrido basado en dos de los algoritmos de aprendizaje no
supervisado (\emph{Unsupervised Learning}) más robustos: K-Means y
DBSCAN. A diferencia del aprendizaje supervisado, estas técnicas no
requieren datos etiquetados previamente, lo cual es idóneo para el
descubrimiento de regímenes de mercado ocultos en series temporales
multidimensionales. A continuación, se detallan sus fundamentos
matemáticos.

\subsubsection{Algoritmo de Agrupamiento
K-Means}\label{algoritmo-de-agrupamiento-k-means}

El algoritmo K-Means es un método de iterativo que busca dividir un
conjunto de \(n\) observaciones en un número predefinido de \(k\)
clústeres disjuntos, de forma que cada observación pertenezca al clúster
con el centroide (media) más cercano. Matemáticamente, dado un conjunto
de vectores de características
\(X = \text{\{}x_{1},x_{2},\ldots,x_{n}\text{\}}\ \) donde cada
\(x_{i} \in \mathbb{R}^{\mathbb{d}}\) (siendo \(d\) el número de
variables o indicadores financieros), el algoritmo tiene como objetivo
minimizar la varianza intra-clúster. Esto se formula como la
minimización de la suma de los errores cuadráticos (\emph{Within-Cluster
Sum of Squares}, WCSS):

\[\arg{\min_{C}{\sum_{i = 1}^{k}{\sum_{x \in C_{i}}^{}\left| \left| x - \mu_{i} \right| \right|^{2}}}}\]

\protect\phantomsection\label{_Toc227174292}{}Ecuación . Minimización de
la suma de los errores cuadráticos

Donde \(C = \text{\{}C_{1},C_{2},\ldots,C_{k}\text{\}}\) representa el
conjunto de clústeres y \(\mu_{i}\) es el centroide del clúster
\(C_{i}\), calculado como la media aritmética de los puntos asignados a
dicho clúster:

\[\mu_{i} = \frac{1}{\left| C_{i} \right|}\sum_{x \in C_{i}}^{}x\]

\protect\phantomsection\label{_Ref227165937}{}Ecuación . Media
aritmética de los puntos del clúster

El proceso de optimización utiliza comúnmente la distancia euclídea como
métrica de similitud en el hiperespacio \(d -\)dimensional:

\[d(x,y) = \sqrt{\sum_{j = 1}^{d}\left( x_{j} - y_{j} \right)^{2}}\]

\protect\phantomsection\label{_Toc227174294}{}Ecuación . Distancia
euclídea

\begin{figure}
\centering
\includegraphics[width=5.76667in,height=4.01906in]{media/image5.PNG}
\caption{\protect\phantomsection\label{_Toc227174289}{}Figura :
Clasificación de activos en el espacio riesgo-rentabilidad mediante el
algoritmo KMeans. Los puntos sólidos representan los centroides µi
calculados según la \hyperref[_Ref227165937]{Ecuación 2}, mientras que
las regiones sombreadas delimitan las fronteras de decisión de cada
clúster. Fuente: Elaboración propia.}
\end{figure}

A pesar de su eficiencia computacional \(\mathcal{O}(nki)\) , el
algoritmo clásico presenta dos limitaciones críticas en finanzas: asume
que los clústeres tienen forma esférica (convergencia isométrica) y, al
carecer de un concepto de densidad, se ve forzado a asignar todos los
puntos a un clúster. Como indica la \hyperref[_Ref227165937]{Ecuación
2}, un valor atípico extremo alterará el sumatorio, desplazando el
centroide \(\mu_{i}\) y sesgando la identificación del régimen de
mercado.

\subsubsection{Filtro de densidad topológica
DBSCAN}\label{filtro-de-densidad-topoluxf3gica-dbscan}

Para mitigar la fragilidad del K-Means ante valores extremos, se
implementa el algoritmo DBSCAN (\emph{Density-Based Spatial Clustering
of Applications with Noise}). A diferencia de los métodos de partición,
DBSCAN define los clústeres como regiones continuas de alta densidad de
puntos separadas por regiones de baja densidad.

La topología de DBSCAN se rige fundamentalmente por dos hiperparámetros
espaciales:

\begin{itemize}
\item
  \(\epsilon\) (\emph{Epsilon}): El radio máximo de vecindad alrededor
  de un punto.
\item
  \(MinPts\): El número mínimo de puntos requeridos dentro de la
  vecindad \(\epsilon\) para considerar que existe una región densa.
\end{itemize}

Para cualquier punto \(p\), su \(\epsilon\) - vecindad se define
matemáticamente como el subconjunto de puntos \(q\) en la base de datos
\(D\) cuya distancia a \(p\) es menor o igual a \(\epsilon\):

\[N_{\epsilon}(p) = \text{\{}q \in D \mid dist(p,q) \leq \epsilon\text{\}}\]

\protect\phantomsection\label{_Toc227174295}{}Ecuación . Vecindad

En base a esta vecindad, el algoritmo clasifica todas las observaciones
del hiperespacio en tres categorías topológicas distintas:

\begin{enumerate}
\def\labelenumi{\arabic{enumi}.}
\item
  Puntos Núcleo (\emph{Core Points}): Un punto \(p\) es un núcleo si
  \(\left| N_{\epsilon}(p) \right| \geq MinPts\). Estos puntos conforman
  el interior (esqueleto) de un régimen de mercado.
\item
  Puntos Frontera (\emph{Border Points}): Un punto \(q\) es frontera si
  \(\left| N_{\epsilon}(q) \right| < MinPts,\) pero pertenece a la
  \(\epsilon\) - vecindad de un punto núcleo.
\item
  Ruido (\emph{Noise} / \emph{Outliers}): Cualquier punto \(o\) que no
  es ni núcleo ni frontera. Matemáticamente, son puntos aislados que no
  alcanzan la densidad crítica.
\end{enumerate}

\begin{figure}
\centering
\includegraphics[width=4.43701in,height=3.72824in]{media/image6.png}
\caption{\protect\phantomsection\label{_Toc227174290}{}Figura :
Clasificación topológica en DBSCAN: Identificación de puntos núcleo,
frontera y ruido en función del radio ϵ y MinPts. Fuente: Elaboración
propia.}
\end{figure}

En el marco de esta investigación, el objetivo primario de DBSCAN no es
la clasificación final de los regímenes, sino la capacidad para detectar
los puntos clasificados como "Ruido". Al asignar matemáticamente la
etiqueta \(- 1\) a las anomalías idiosincrásicas (como el \emph{Flash
Crash} previamente ilustrado), el algoritmo permite purgar el
hiperespacio de características antes de someterlo a la minimización de
la varianza del K-Means.

\subsection{Estado del arte y justificación de la arquitectura
propuesta}\label{estado-del-arte-y-justificaciuxf3n-de-la-arquitectura-propuesta}

La aplicación de técnicas de \emph{Machine Learning} no supervisado en
la optimización de carteras ha experimentado un crecimiento exponencial.
Sin embargo, la comunidad científica aún debate cómo gestionar la
disrupción que generan las anomalías del mercado sobre los algoritmos
clásicos. A continuación, se articula el estado de la cuestión siguiendo
la evolución desde el problema tradicional hasta la propuesta
arquitectónica de esta investigación.

\subsubsection{El problema del Momentum tradicional en entornos
volátiles}\label{el-problema-del-momentum-tradicional-en-entornos-voluxe1tiles}

La selección de activos para conformar carteras de \emph{momentum} ha
incrementado su complejidad metodológica. Según advierte (Aslam, 2025),
esta toma de decisiones representa en la actualidad un desafío mayúsculo
debido a la alta volatilidad inherente a los mercados. El autor sugiere
que el \emph{momentum} transversal estándar es a menudo insuficiente en
estos entornos de estrés, lo que justifica de manera imperativa la
aplicación de algoritmos de \emph{Machine Learning} para el desarrollo
de una "estrategia de \emph{momentum} gestionada por riesgo".

En paralelo, (Shiraj, y otros, 2024) exponen en su análisis que los
mercados financieros están expuestos de forma recurrente a movimientos
extremos de precios repentinos. Estos eventos anómalos o \emph{crashes},
si no son detectados y aislados a tiempo, actúan como un factor de
contaminación estadística capaz de invalidar por completo los modelos
predictivos y de clasificación subyacentes.

\subsubsection{El intento de solución mediante K-Means y sus
limitaciones}\label{el-intento-de-soluciuxf3n-mediante-k-means-y-sus-limitaciones}

Como respuesta a esta necesidad de adaptación dinámica, la literatura
reciente ha explorado modelos de agrupamiento basados en partición.
(Aslam, 2025) propone la utilización del algoritmo K-Means para
clasificar el universo de activos basándose en características
multidimensionales de riesgo y retorno. Su enfoque permite la
identificación algorítmica de clústeres compuestos por las denominadas
"acciones optimistas", superando los resultados de los deciles
tradicionales.

Adicionalmente, otros trabajos contemporáneos como el de (Luan \& Hamp,
2025) evidencian la necesidad de automatizar la clasificación de
regímenes en series temporales multidimensionales, proponiendo variantes
avanzadas del K-Means (como el \emph{Sliced Wasserstein K-Means}) para
tratar de lidiar con la complejidad de las distribuciones financieras.

No obstante, esta línea de solución presenta una limitación técnica
crítica advertida por (Nyffeler, 2024): el algoritmo K-Means tradicional
posee una sensibilidad matemática inherente a los valores atípicos
(\emph{outliers}). Dado que minimiza la suma de los errores cuadráticos
global, la presencia de anomalías genera una severa inestabilidad en la
asignación de los clústeres y una distorsión artificial en la ubicación
de los centroides. Esta vulnerabilidad es ejemplificada de manera
concluyente por (Bin, 2020), quien documenta cómo la caída extrema de
una sola empresa (Apache Corp, -70\%) es capaz de desplazar los
centroides y comprometer el rendimiento de toda una cartera K-Means.

\subsubsection{Propuesta de Solución: Arquitectura Híbrida y Marco de
Evaluación}\label{propuesta-de-soluciuxf3n-arquitectura-huxedbrida-y-marco-de-evaluaciuxf3n}

La revisión del estado del arte evidencia una brecha: los algoritmos
particionales (K-Means) agrupan eficientemente el riesgo, pero son
incapaces de discriminar eventos irracionales, mientras que los
algoritmos de densidad (DBSCAN) identifican las anomalías, pero no
segmentan el mercado direccionalmente.

Para resolver esta problemática, el presente Trabajo propone una
solución híbrida. Integrando las advertencias de (Shiraj, y otros, 2024)
y la necesidad de gestión de riesgo de (Aslam, 2025), se establece una
arquitectura secuencial: el algoritmo DBSCAN actuará como preprocesador
para detectar y aislar matemáticamente las anomalías, protegiendo así al
modelo. Posteriormente, K-Means clasificará el espacio de datos ya
purgado, asegurando que los centroides representen los verdaderos
regímenes de "acciones optimistas" sin la distorsión del ruido
estructural.

Para validar de forma empírica la superioridad de esta arquitectura y
mitigar el riesgo de sobreajuste (\emph{overfitting}) advertido por
(López de Prado, 2018), los hiperparámetros se optimizarán en una
ventana de entrenamiento (\emph{in-sample}) y se evaluarán en un periodo
cronológico estrictamente posterior (\emph{out-of-sample}). La
investigación contrastará tres escenarios de inversión:

\begin{enumerate}
\def\labelenumi{\arabic{enumi}.}
\item
  Benchmark (Referencia): Rentabilidad pasiva del índice S\&P 500
  (Comprar y mantener).
\item
  Modelo Base: Estrategia algorítmica utilizando K-Means aislado.
\item
  Modelo Híbrido (Propuesta): Implementación secuencial de DBSCAN
  seguido de K-Means.
\end{enumerate}

La robustez de los modelos se medirá a través de métricas
institucionales clave: el \emph{Ratio de Sharpe} (rendimiento ajustado
al riesgo) y el \emph{Maximum Drawdown} (riesgo de caída máxima). El
objetivo final de la experimentación será demostrar matemáticamente que
el filtrado previo de anomalías reduce de forma efectiva las pérdidas de
la cartera durante los escenarios de estrés del mercado.

\section{Objetivos concretos y metodología de
trabajo}\label{objetivos-concretos-y-metodologuxeda-de-trabajo}

En este capítulo se definen los objetivos del trabajo y la metodología
que se seguirá para su desarrollo. Este apartado actúa como nexo entre
el estudio del contexto y estado del arte presentado anteriormente y la
implementación práctica de la propuesta planteada.

El trabajo se enmarca en el ámbito del análisis de datos financieros
mediante técnicas de \emph{Machine Learning}, con el objetivo de
desarrollar y evaluar una arquitectura algorítmica híbrida capaz de
identificar patrones de comportamiento en activos financieros y reducir
el impacto de anomalías en entornos de alta volatilidad.

En primer lugar, se establece el objetivo general del trabajo, que
define el propósito principal de la investigación. A continuación, se
desglosa dicho objetivo en una serie de objetivos específicos que
permiten estructurar el desarrollo del proyecto en tareas concretas.
Finalmente, se describe la metodología de trabajo, detallando las fases,
técnicas y herramientas que se emplearán para alcanzar los objetivos
definidos.

\subsection{Objetivo general}\label{objetivo-general}

El objetivo general de este trabajo es diseñar, implementar y evaluar
una arquitectura híbrida basada en técnicas de \emph{clustering} (DBSCAN
y K-Means) que permita identificar patrones de comportamiento en activos
financieros y mejorar la robustez de las estrategias de \emph{momentum}
mediante el filtrado previo de anomalías.

Para ello, se desarrollará un sistema capaz de agrupar activos con
dinámicas similares tras un proceso de depuración del ruido estructural,
generando señales indicativas de posibles oportunidades de inversión o
eventos atípicos.

El sistema propuesto será evaluado con el fin de analizar su utilidad en
la interpretación del comportamiento del mercado y en la toma de
decisiones basada en datos.

\subsection{Objetivos específicos}\label{objetivos-especuxedficos}

Para alcanzar el objetivo general, se plantean los siguientes objetivos
específicos:

\begin{enumerate}
\def\labelenumi{\arabic{enumi}.}
\item
  \textbf{Recopilar datos históricos} de precios de activos financieros
  a partir de fuentes públicas.
\item
  \textbf{Analizar la estructura y características de los datos}
  obtenidos, incluyendo variables como precios de apertura, cierre,
  máximos, mínimos y volumen.
\item
  \textbf{Preprocesar y limpiar el conjunto de datos}, tratando valores
  nulos y datos inconsistentes.
\item
  \textbf{Transformar las series temporales en variables relevantes},
  como retornos o medidas de volatilidad.
\item
  \textbf{Implementar una arquitectura secuencial} basada en DBSCAN y
  K-Means para el análisis del comportamiento de los activos.
\item
  \textbf{Aplicar DBSCAN} como técnica de filtrado para identificar y
  eliminar observaciones anómalas (outliers).
\item
  \textbf{Aplicar K-Means} sobre el conjunto de datos depurado para
  identificar clústeres representativos de regímenes de mercado.
\item
  \textbf{Analizar los clusters obtenidos e interpretar su significado}
  en términos financieros.
\item
  \textbf{Integrar una estrategia de momentum} como marco de evaluación
  del sistema propuesto.
\item
  \textbf{Comparar el rendimiento} del modelo híbrido frente a un modelo
  basado únicamente en K-Means.
\item
  \textbf{Evaluar el rendimiento del sistema} utilizando datos no vistos
  previamente.
\item
  \textbf{Validar la utilidad del modelo} en un contexto experimental.
\end{enumerate}

\subsection{Metodología del trabajo}\label{metodologuxeda-del-trabajo}

La metodología del presente trabajo se estructura como un proceso
iterativo de análisis, modelado y evaluación de datos financieros, con
el objetivo de desarrollar un sistema capaz de identificar patrones de
comportamiento en activos y evaluar el impacto del filtrado de anomalías
en la robustez de estrategias de inversión basadas en \emph{momentum}.

\subsubsection{Obtención de datos}\label{obtenciuxf3n-de-datos}

En una primera fase, se realiza la obtención de datos históricos de
mercados financieros a partir de fuentes públicas en línea. Estos datos
incluyen información temporal de los activos, como precios de apertura,
cierre, máximos, mínimos y volumen de negociación.

La elección de este tipo de datos se fundamenta en su uso habitual
dentro del análisis cuantitativo de mercados, así como en su capacidad
para reflejar dinámicas reales de comportamiento en los precios.

\subsubsection{Procesamiento y preparación de
datos}\label{procesamiento-y-preparaciuxf3n-de-datos}

Una vez recopilados los datos, se lleva a cabo un proceso de preparación
cuyo objetivo es garantizar su calidad y adecuación para los modelos de
aprendizaje automático.

En esta fase se realiza la limpieza del conjunto de datos, incluyendo la
gestión de valores nulos, la corrección de inconsistencias y la
eliminación de registros incompletos. Posteriormente, se aplican
transformaciones sobre las series temporales con el fin de generar
variables más informativas, como retornos porcentuales o medidas de
volatilidad.

Este proceso resulta fundamental para adaptar los datos en bruto a un
formato adecuado para su posterior análisis.

\subsubsection{Análisis exploratorio de datos
(EDA)}\label{anuxe1lisis-exploratorio-de-datos-eda}

A continuación, se realiza un análisis exploratorio de los datos con el
objetivo de comprender su estructura interna, identificar patrones
iniciales y detectar posibles anomalías o comportamientos relevantes.

Este análisis incluye el estudio de la evolución temporal de los
activos, la relación entre variables y la comparación de comportamientos
entre distintos activos financieros. Los resultados obtenidos en esta
fase permiten orientar la selección de los modelos de \emph{clustering}
y la definición de la estrategia de análisis posterior.

\subsubsection{Implementación de la arquitectura
híbrida}\label{implementaciuxf3n-de-la-arquitectura-huxedbrida}

Sobre los datos ya procesados, se implementa una arquitectura secuencial
en dos etapas:

En primer lugar, se aplica el algoritmo DBSCAN como mecanismo de
filtrado topológico, con el objetivo de identificar observaciones
anómalas y clasificarlas como ruido. Posteriormente, se elimina dicho
ruido del conjunto de datos y se aplica el algoritmo K-Means sobre el
espacio depurado, con el fin de identificar clústeres representativos de
distintos regímenes de mercado.

Esta aproximación permite reducir la influencia de valores atípicos en
la estimación de los centroides, mejorando la estabilidad y coherencia
de los clústeres obtenidos.

\subsubsection{Estrategia de momentum y detección de
anomalías}\label{estrategia-de-momentum-y-detecciuxf3n-de-anomaluxedas}

A partir de los \emph{clusters} obtenidos, se desarrolla una estrategia
de \emph{momentum} que actúa como marco de evaluación del sistema
propuesto.

En particular, se analizan transiciones entre distintos \emph{clusters}
y desviaciones significativas respecto a patrones establecidos, lo que
permite identificar posibles señales de inversión o situaciones anómalas
en el mercado.

\subsubsection{Evaluación del sistema}\label{evaluaciuxf3n-del-sistema}

Finalmente, el sistema se evalúa utilizando datos históricos no
empleados durante la fase de construcción de los modelos, con el
objetivo de analizar su capacidad de generalización.

La evaluación se centra en la coherencia de las señales generadas, así
como en la estabilidad del modelo en distintos escenarios de mercado.

\subsubsection{Entorno de desarrollo y
herramientas}\label{entorno-de-desarrollo-y-herramientas}

El desarrollo del sistema se lleva a cabo en entorno \emph{Python},
utilizando \emph{JupyterLab} como entorno principal de trabajo, junto
con bibliotecas especializadas en análisis de datos y \emph{Machine
Learning} como \emph{pandas}, \emph{numpy} y \emph{scikit-learn}.

\section{Referencias
bibliográficas}\label{referencias-bibliogruxe1ficas}

Aslam, B. (2025). Identifying optimistic stocks with K-means clustering
algorithm. \emph{International Review of Economics \& Finance}, 104579.

Bin, S. (2020). K-means stock clustering analysis based on historical
price movements and financial ratios (Senior thesis). \emph{Claremont
McKenna College}.

Daniel, K. M. (2016). Momentum crashes. \emph{Journal of Financial
Economics}, 221-247.

Jegadeesh, N. T. (1993). Returns to Buying Winners and Selling Losers:
Implications for Stock Market Efficiency. \emph{The Journal of Finance},
65--91.

López de Prado, M. (2018). \emph{Advances in Financial Machine
Learning.} Wiley.

Luan, Q., \& Hamp, J. (2025). Automated regime classification in
multidimensional time series data using sliced Wasserstein k-means
clustering. \emph{Data Science in Finance and Economics}, 387--418.

Nyffeler, L. (2024). K-means clustering with Mahalanobis distance for
stock classification into value and growth portfolios (Bachelor's
thesis). \emph{Zurich University of Applied Sciences (ZHAW)}.

Shiraj, M., Rahman, M., Al-Imran, M., Liza, M., Murshed, M., \& Akhter,
N. (2024). Anomaly detection in financial time series data via mapper
algorithm and DBSCAN clustering. \emph{World Journal of Advanced
Engineering Technology and Sciences}, 70-84.

\end{document}
